% Medidas ambientais e sociais ao longo do ciclo da ação
\begin{tikzpicture}[node distance=0.08cm,
    medida/.style={draw=azulpsh!60, fill=white, rounded corners=2pt, align=left,
                   text width=3.35cm, inner sep=4pt, font=\scriptsize,
                   minimum height=4.3cm, anchor=north}]
  \node[fase=3.3cm] (f1) {Planejamento e\\ seleção};
  \node[fase=3.3cm, right=of f1] (f2) {Projeto e\\ licenciamento};
  \node[fase=3.3cm, right=of f2] (f3) {Obras e\\ instalação};
  \node[fase=3.3cm, fill=verdeclaro, right=of f3] (f4) {Operação\\ e entrega};

  \node[medida, below=0.3cm of f1.south] {%
    \textbullet\ consultas e métodos participativos na seleção\\
    \textbullet\ busca ativa de famílias em áreas remotas\\
    \textbullet\ inclusão de grupos vulneráveis e comunidades tradicionais\\
    \textbullet\ verificação da titularidade das áreas, com anuências e servidões};
  \node[medida, below=0.3cm of f2.south] {%
    \textbullet\ estudos ambientais e sociais\\
    \textbullet\ outorga do uso da água e licenças\\
    \textbullet\ normas técnicas de perfuração\\
    \textbullet\ dimensionamento da solução de esgoto para cada domicílio};
  \node[medida, below=0.3cm of f3.south] {%
    \textbullet\ saúde e segurança do trabalho, com EPI\\
    \textbullet\ sinalização e proteção de valas\\
    \textbullet\ gestão de resíduos, ruído e poeira\\
    \textbullet\ testes de estanqueidade após a instalação};
  \node[medida, below=0.3cm of f4.south] {%
    \textbullet\ capacitação em uso e manutenção\\
    \textbullet\ modelo de gestão e tarifa social\\
    \textbullet\ manejo adequado do lodo\\
    \textbullet\ suporte técnico e avaliação semestral};

  \node[transversal=14.3cm, below=4.85cm of $(f2.south)!0.5!(f3.south)$, anchor=north] {%
    \textbf{Em todas as fases:} reuniões comunitárias em linguagem simples; canal de
    queixas e reclamações divulgado nas comunidades, com tempo de resposta monitorado a
    cada semestre.};
\end{tikzpicture}
